\section{Background and Related Works}

\noindent\textbf{MAD preliminary.} In MAD, LLM agents iteratively exchange progress. Suppose $N$ participating agents. Let $C_{t}^i$ be the output (i.e., reasoning and answer) from agent $i$ at turn $t$. Then, in all-to-all MAD, agent $i$ generates the next-turn output, $C_{t+1}^i$, based on the following context:

[Question] + [turn $1,\dots,t-1$ history] + $C_t^i$ + $DebateFormat(C_{t}^1,\dots,C_t^{i-1}, C_t^{i+1}\dots C_t^N)$

where $DebateFormat$ is a string manipulation function that concatenates peer outputs and adds a prompting hint to guide agent $i$ to reflect on peer contexts, such as \textit{"These are the solutions from other agents [Insert peer solutions here]...Using the solutions from other agents as additional information, provide your answer..."}~\citep{Du2024Improving}.

While MAD improves LLM-agent accuracy, it scales poorly with more agents due to rapid context growth. Under all-to-all communication, each agent’s context grows as $\mathcal{O}(N^2t)$, increasing computation and KV-cache pressure while also risking long-context degradation from diluted salient information~\citep{Liu2024Lost, Laban2026Llm}.

\noindent\textbf{Efficient MAD solutions.} Existing works leverage pre-debate signals to improve MAD. One line of work uses agent-level reliability signals to identify agents whose answers already appear reliable and can therefore conclude without any debate~\citep{Chen2025Sid, Eo2025Debate}. Another line uses cross-agent similarity signals to remove communication links that appear redundant~\citep{Zeng2025S2Mad, Liu2025Groupdebate, Zhu2026Demystifying}.
